\section{光绪新政前后的账本叙事：地方社会与制度重组}
1875年至1908年的改革、灾荒、边疆、战争和公共信息网络在账本中彼此牵动。本节保留实录、部院档、海关及财政材料、地方档、报刊和中外外交文本各自的观察位置；设局、立法、购舰、设省、办学或宣言都被视为过程的一个节点，而非已经覆盖全境的结果。
\subsection{1870—1879：事件、文书与地方后果}
这一组节点按账本的日期顺序排列。相邻事件并不必然构成单一因果，却共同显示信息、资源与地方行动如何在同一时段交错。
\hypertarget{ledger:EQ-1875-GUANGXU-ACCESSION}{}\paragraph{光绪元年（1875年）}光绪帝即位，慈禧、慈安继续掌握关键政治影响。核心取证为《清德宗实录》光绪元年相关条；宫中朱批奏折、内阁大库题本与《清史稿》相关志传。这类上谕、题本或部议首先证明机构处理了何种命令、呈报与责任，而不自动证明地方收到、执行或接受。账本所列条件为：幼主继承使摄政与名分安排再度出现。其后果被限定为：清末改革和对外危机将在这一权力格局下发展。取证保留：以光绪元年（1875年）的同期文书为定位起点；官修记录、奏报或外交文本服务特定机构立场，具体日期、规模、地方执行与对手视角须以独立材料复核。
\hypertarget{ledger:EQ-1875-YUNNAN-MUSLIM-WAR}{}\paragraph{光绪元年（1875年）}杜文秀政权覆亡前后，云南回民战争进入惨烈收束阶段。核心取证为《清德宗实录》光绪元年相关条；军机处档、战事方略及地方奏报。编年记录适于钉住事件在朝廷文书中的时间位置，却需以地方、对方或物质材料检验其范围和后果。账本所列条件为：地方军政、矿业贸易与族群冲突长期累积。其后果被限定为：死亡、迁徙与复垦规模的估计须保持不确定性。取证保留：以光绪元年（1875年）的同期文书为定位起点；官修记录、奏报或外交文本服务特定机构立场，具体日期、规模、地方执行与对手视角须以独立材料复核。
\hypertarget{ledger:EQ-1876-CHEFOO-CONVENTION}{}\paragraph{光绪二年（1876年）}《烟台条约》签订，扩大英国在华通商和内地活动权利。核心取证为《清德宗实录》光绪二年相关条；《筹办夷务始末》相关年卷与中外照会、条约文本。这类上谕、题本或部议首先证明机构处理了何种命令、呈报与责任，而不自动证明地方收到、执行或接受。账本所列条件为：马嘉理事件与英国外交压力推动谈判。其后果被限定为：条约权利的地区落实和地方反应需要逐案追踪。取证保留：以光绪二年（1876年）的同期文书为定位起点；官修记录、奏报或外交文本服务特定机构立场，具体日期、规模、地方执行与对手视角须以独立材料复核。
\hypertarget{ledger:EQ-1876-NORTH-CHINA-FAMINE}{}\paragraph{光绪二年（1876年）}丁戊奇荒在华北及西北造成广泛饥荒、迁徙和救济行动。核心取证为《清德宗实录》光绪二年相关条；地方志、督抚奏折及地方档案。编年记录适于钉住事件在朝廷文书中的时间位置，却需以地方、对方或物质材料检验其范围和后果。账本所列条件为：旱灾、粮价、交通和行政能力共同塑造灾情。其后果被限定为：伤亡总数争议极大，应按地区、材料与估算口径说明。取证保留：以光绪二年（1876年）的同期文书为定位起点；官修记录、奏报或外交文本服务特定机构立场，具体日期、规模、地方执行与对手视角须以独立材料复核。
\hypertarget{ledger:EQ-1877-FAMINE-RELIEF-NETWORKS}{}\paragraph{光绪三年（1877年）}赈济网络在灾区展开，官办、绅办和传教团体的救助并存。核心取证为《清德宗实录》光绪三年相关条；地方志、督抚奏折及地方档案。编年记录适于钉住事件在朝廷文书中的时间位置，却需以地方、对方或物质材料检验其范围和后果。账本所列条件为：灾害规模超出常规仓储和地方财政能力。其后果被限定为：救助记录反映组织覆盖，不等于所有灾民获救。取证保留：以光绪三年（1877年）的同期文书为定位起点；官修记录、奏报或外交文本服务特定机构立场，具体日期、规模、地方执行与对手视角须以独立材料复核。
\hypertarget{ledger:EQ-1878-Zuo-ZONGTANG-XINJIANG}{}\paragraph{光绪四年（1878年）}左宗棠军收复新疆大部，阿古柏政权崩解。核心取证为《清德宗实录》光绪四年相关条；军机处档、战事方略及地方奏报。编年记录适于钉住事件在朝廷文书中的时间位置，却需以地方、对方或物质材料检验其范围和后果。账本所列条件为：清廷集中湘军、饷源和西北运输资源进行反攻。其后果被限定为：军事收复与后续行政、绿洲社会重建必须分开叙述。取证保留：以光绪四年（1878年）的同期文书为定位起点；官修记录、奏报或外交文本服务特定机构立场，具体日期、规模、地方执行与对手视角须以独立材料复核。
\hypertarget{ledger:EQ-1878-XINJIANG-PROVINCE-DEBATE}{}\paragraph{光绪四年（1878年）}朝廷讨论新疆设省、军政与财政安排。核心取证为《清德宗实录》光绪四年相关条；宫中朱批奏折、内阁大库题本与《清史稿》相关志传。这类上谕、题本或部议首先证明机构处理了何种命令、呈报与责任，而不自动证明地方收到、执行或接受。账本所列条件为：战后长期驻防成本和俄国边境压力要求重建制度。其后果被限定为：设省方案反映中央目标，地方实施尚需持续资源。取证保留：以光绪四年（1878年）的同期文书为定位起点；官修记录、奏报或外交文本服务特定机构立场，具体日期、规模、地方执行与对手视角须以独立材料复核。
\hypertarget{ledger:EQ-1879-RYUKYU-ANNEXATION}{}\paragraph{光绪五年（1879年）}日本吞并琉球，清廷的交涉未能恢复旧有册封关系。核心取证为《清德宗实录》光绪五年相关条；《筹办夷务始末》相关年卷与中外照会、条约文本。这类上谕、题本或部议首先证明机构处理了何种命令、呈报与责任，而不自动证明地方收到、执行或接受。账本所列条件为：日本国家建构与东亚秩序转变相互作用。其后果被限定为：清方外交文书与琉球、日本材料对权利表述不同。取证保留：以光绪五年（1879年）的同期文书为定位起点；官修记录、奏报或外交文本服务特定机构立场，具体日期、规模、地方执行与对手视角须以独立材料复核。
\subsection{1880—1889：事件、文书与地方后果}
这一组节点按账本的日期顺序排列。相邻事件并不必然构成单一因果，却共同显示信息、资源与地方行动如何在同一时段交错。
\hypertarget{ledger:EQ-1880-ILI-NEGOTIATIONS}{}\paragraph{光绪六年（1880年）}清俄围绕伊犁归还和边界条件持续谈判。核心取证为《清德宗实录》光绪六年相关条；《筹办夷务始末》相关年卷与中外照会、条约文本。这类上谕、题本或部议首先证明机构处理了何种命令、呈报与责任，而不自动证明地方收到、执行或接受。账本所列条件为：俄军占据伊犁后，清廷在军事与财政限制下寻求外交解决。其后果被限定为：边界线、赔款和贸易条款须逐项区分。取证保留：以光绪六年（1880年）的同期文书为定位起点；官修记录、奏报或外交文本服务特定机构立场，具体日期、规模、地方执行与对手视角须以独立材料复核。
\hypertarget{ledger:EQ-1881-SAINT-PETERSBURG-TREATY}{}\paragraph{光绪七年（1881年）}《中俄伊犁条约》签订，伊犁大部归还清廷并附带领土、赔款和贸易条件。核心取证为《清德宗实录》光绪七年相关条；《筹办夷务始末》相关年卷与中外照会、条约文本。这类上谕、题本或部议首先证明机构处理了何种命令、呈报与责任，而不自动证明地方收到、执行或接受。账本所列条件为：长期谈判与西北军事实力恢复共同作用。其后果被限定为：条约地图与实际边界管理不是同一问题。取证保留：以光绪七年（1881年）的同期文书为定位起点；官修记录、奏报或外交文本服务特定机构立场，具体日期、规模、地方执行与对手视角须以独立材料复核。
\hypertarget{ledger:EQ-1882-IMO-INCIDENT}{}\paragraph{光绪八年（1882年）}朝鲜壬午兵变后清军入朝，日本也扩大影响。核心取证为《清德宗实录》光绪八年相关条；《筹办夷务始末》相关年卷与中外照会、条约文本。这类上谕、题本或部议首先证明机构处理了何种命令、呈报与责任，而不自动证明地方收到、执行或接受。账本所列条件为：朝鲜内政、军制改革与列强竞争交织。其后果被限定为：清、朝、日文献中的“保护”与控制含义不同。取证保留：以光绪八年（1882年）的同期文书为定位起点；官修记录、奏报或外交文本服务特定机构立场，具体日期、规模、地方执行与对手视角须以独立材料复核。
\hypertarget{ledger:EQ-1883-SINO-FRENCH-WAR-BEGIN}{}\paragraph{光绪九年（1883年）}中法因越南北部和宗藩关系冲突升级为战争。核心取证为《清德宗实录》光绪九年相关条；军机处档、战事方略及地方奏报。编年记录适于钉住事件在朝廷文书中的时间位置，却需以地方、对方或物质材料检验其范围和后果。账本所列条件为：法国殖民扩张与清廷在越南的传统关系碰撞。其后果被限定为：战事跨越越南、台湾和华南，不能只从北京外交看。取证保留：以光绪九年（1883年）的同期文书为定位起点；官修记录、奏报或外交文本服务特定机构立场，具体日期、规模、地方执行与对手视角须以独立材料复核。
\hypertarget{ledger:EQ-1884-FUZHOU-NAVAL-BATTLE}{}\paragraph{光绪十年（1884年）}马江海战中福建水师遭法国舰队重创。核心取证为《清德宗实录》光绪十年相关条；军机处档、战事方略及地方奏报。编年记录适于钉住事件在朝廷文书中的时间位置，却需以地方、对方或物质材料检验其范围和后果。账本所列条件为：清军舰队训练、指挥与武器体系面临现代海战挑战。其后果被限定为：伤亡和舰只损失应依据中法记录和船政资料核验。取证保留：以光绪十年（1884年）的同期文书为定位起点；官修记录、奏报或外交文本服务特定机构立场，具体日期、规模、地方执行与对手视角须以独立材料复核。
\hypertarget{ledger:EQ-1884-TAIWAN-BLOCKADE}{}\paragraph{光绪十年（1884年）}法军进攻台湾北部并封锁，台湾成为中法战争重要战场。核心取证为《清德宗实录》光绪十年相关条；军机处档、战事方略及地方奏报。编年记录适于钉住事件在朝廷文书中的时间位置，却需以地方、对方或物质材料检验其范围和后果。账本所列条件为：台湾港口、煤矿和海峡航路具有战略价值。其后果被限定为：地方抵抗、商业中断与原住民地区情况需分区考察。取证保留：以光绪十年（1884年）的同期文书为定位起点；官修记录、奏报或外交文本服务特定机构立场，具体日期、规模、地方执行与对手视角须以独立材料复核。
\hypertarget{ledger:EQ-1885-TIANJIN-CHINA-FRANCE-TREATY}{}\paragraph{光绪十一年（1885年）}《中法新约》确认越南脱离清朝宗藩体系，战争结束。核心取证为《清德宗实录》光绪十一年相关条；《筹办夷务始末》相关年卷与中外照会、条约文本。这类上谕、题本或部议首先证明机构处理了何种命令、呈报与责任，而不自动证明地方收到、执行或接受。账本所列条件为：军事与财政压力推动双方议和。其后果被限定为：条约确认与边界勘界、华南地方影响应分别处理。取证保留：以光绪十一年（1885年）的同期文书为定位起点；官修记录、奏报或外交文本服务特定机构立场，具体日期、规模、地方执行与对手视角须以独立材料复核。
\hypertarget{ledger:EQ-1885-TAIWAN-PROVINCE}{}\paragraph{光绪十一年（1885年）}清廷决定将台湾建省，加强行政、军防和基础设施经营。核心取证为《清德宗实录》光绪十一年相关条；宫中朱批奏折、内阁大库题本与《清史稿》相关志传。这类上谕、题本或部议首先证明机构处理了何种命令、呈报与责任，而不自动证明地方收到、执行或接受。账本所列条件为：中法战争显示台湾海防与行政地位的重要性。其后果被限定为：建省不等于全岛被均质治理，东部和原住民地区尤需谨慎。取证保留：以光绪十一年（1885年）的同期文书为定位起点；官修记录、奏报或外交文本服务特定机构立场，具体日期、规模、地方执行与对手视角须以独立材料复核。
\hypertarget{ledger:EQ-1886-BEIYANG-NAVY-EXPANSION}{}\paragraph{光绪十二年（1886年）}北洋海军和沿海防务建设持续推进。核心取证为《清德宗实录》光绪十二年相关条；军机处档、战事方略及地方奏报。编年记录适于钉住事件在朝廷文书中的时间位置，却需以地方、对方或物质材料检验其范围和后果。账本所列条件为：中法战争和日本扩张促使海军现代化。其后果被限定为：购舰数量不等于训练、维护、弹药和指挥能力。取证保留：以光绪十二年（1886年）的同期文书为定位起点；官修记录、奏报或外交文本服务特定机构立场，具体日期、规模、地方执行与对手视角须以独立材料复核。
\hypertarget{ledger:EQ-1887-TAIWAN-RAILWAY-PLANNING}{}\paragraph{光绪十三年（1887年）}台湾巡抚刘铭传推动铁路、电报、矿务和新式行政项目。核心取证为《清德宗实录》光绪十三年相关条；户部奏销、盐政、河工或海关档案。这类上谕、题本或部议首先证明机构处理了何种命令、呈报与责任，而不自动证明地方收到、执行或接受。账本所列条件为：建省后需要连接港口、城市和军防节点。其后果被限定为：项目经费、路段使用与社会影响需具体核验。取证保留：以光绪十三年（1887年）的同期文书为定位起点；官修记录、奏报或外交文本服务特定机构立场，具体日期、规模、地方执行与对手视角须以独立材料复核。
\hypertarget{ledger:EQ-1888-NAVAL-AND-ARMY-REFORM}{}\paragraph{光绪十四年（1888年）}清廷讨论海军、陆军训练和军工体系的经费与制度问题。核心取证为《清德宗实录》光绪十四年相关条；宫中朱批奏折、内阁大库题本与《清史稿》相关志传。这类上谕、题本或部议首先证明机构处理了何种命令、呈报与责任，而不自动证明地方收到、执行或接受。账本所列条件为：列强战争压力暴露既有军制不足。其后果被限定为：改革奏议与实际编制、操练之间常有距离。取证保留：以光绪十四年（1888年）的同期文书为定位起点；官修记录、奏报或外交文本服务特定机构立场，具体日期、规模、地方执行与对手视角须以独立材料复核。
\hypertarget{ledger:EQ-1889-GUANGXU-MARRIAGE}{}\paragraph{光绪十五年（1889年）}光绪帝大婚并开始形式上的亲政安排。核心取证为《清德宗实录》光绪十五年相关条；宫中朱批奏折、内阁大库题本与《清史稿》相关志传。这类上谕、题本或部议首先证明机构处理了何种命令、呈报与责任，而不自动证明地方收到、执行或接受。账本所列条件为：皇帝成年要求完成宫廷礼制程序。其后果被限定为：慈禧仍对重要决策保有深刻影响。取证保留：以光绪十五年（1889年）的同期文书为定位起点；官修记录、奏报或外交文本服务特定机构立场，具体日期、规模、地方执行与对手视角须以独立材料复核。
\subsection{1890—1899：事件、文书与地方后果}
这一组节点按账本的日期顺序排列。相邻事件并不必然构成单一因果，却共同显示信息、资源与地方行动如何在同一时段交错。
\hypertarget{ledger:EQ-1890-KOREAN-POLITICAL-CONFLICT}{}\paragraph{光绪十六年（1890年）}朝鲜内政纷争持续牵动清日两国在半岛的影响力竞争。核心取证为《清德宗实录》光绪十六年相关条；《筹办夷务始末》相关年卷与中外照会、条约文本。这类上谕、题本或部议首先证明机构处理了何种命令、呈报与责任，而不自动证明地方收到、执行或接受。账本所列条件为：壬午事变后清日均扩大驻军与交涉。其后果被限定为：朝鲜行动者不是外部竞争的被动对象，应纳入叙述。取证保留：以光绪十六年（1890年）的同期文书为定位起点；官修记录、奏报或外交文本服务特定机构立场，具体日期、规模、地方执行与对手视角须以独立材料复核。
\hypertarget{ledger:EQ-1891-MISSIONARY-CONFLICTS}{}\paragraph{光绪十七年（1891年）}各地教案和宗教纠纷继续引发地方治理与外交压力。核心取证为《清德宗实录》光绪十七年相关条；地方志、督抚奏折及地方档案。编年记录适于钉住事件在朝廷文书中的时间位置，却需以地方、对方或物质材料检验其范围和后果。账本所列条件为：土地、诉讼、保护权和谣言与宗教身份交织。其后果被限定为：“教案”包含不同类型纠纷，不能单一化。取证保留：以光绪十七年（1891年）的同期文书为定位起点；官修记录、奏报或外交文本服务特定机构立场，具体日期、规模、地方执行与对手视角须以独立材料复核。
\hypertarget{ledger:EQ-1892-RAILWAY-DEBATE}{}\paragraph{光绪十八年（1892年）}官员、商人和士绅围绕铁路修筑、主权与资金展开争论。核心取证为《清德宗实录》光绪十八年相关条；户部奏销、盐政、河工或海关档案。这类上谕、题本或部议首先证明机构处理了何种命令、呈报与责任，而不自动证明地方收到、执行或接受。账本所列条件为：交通技术、外债和地方利益引发新问题。其后果被限定为：赞成或反对文本不等于实际线路建设和公众意见。取证保留：以光绪十八年（1892年）的同期文书为定位起点；官修记录、奏报或外交文本服务特定机构立场，具体日期、规模、地方执行与对手视角须以独立材料复核。
\hypertarget{ledger:EQ-1893-KOREAN-REFORM-CRISIS}{}\paragraph{光绪十九年（1893年）}朝鲜改革与政治危机使清日关系持续紧张。核心取证为《清德宗实录》光绪十九年相关条；《筹办夷务始末》相关年卷与中外照会、条约文本。这类上谕、题本或部议首先证明机构处理了何种命令、呈报与责任，而不自动证明地方收到、执行或接受。账本所列条件为：半岛内政和驻军安排缺乏稳定共识。其后果被限定为：清日文书各自服务战略目的，需与朝鲜材料互读。取证保留：以光绪十九年（1893年）的同期文书为定位起点；官修记录、奏报或外交文本服务特定机构立场，具体日期、规模、地方执行与对手视角须以独立材料复核。
\hypertarget{ledger:EQ-1894-DONGHAK-PEASANT-WAR}{}\paragraph{光绪二十年（1894年）}东学农民战争促使清日两国派兵入朝，危机迅速国际化。核心取证为《清德宗实录》光绪二十年相关条；军机处档、战事方略及地方奏报。编年记录适于钉住事件在朝廷文书中的时间位置，却需以地方、对方或物质材料检验其范围和后果。账本所列条件为：朝鲜社会改革诉求与政府镇压失败交织。其后果被限定为：农民运动不应只作为清日战争的前奏叙述。取证保留：以光绪二十年（1894年）的同期文书为定位起点；官修记录、奏报或外交文本服务特定机构立场，具体日期、规模、地方执行与对手视角须以独立材料复核。
\hypertarget{ledger:EQ-1894-FIRST-SINO-JAPANESE-WAR}{}\paragraph{光绪二十年（1894年）}甲午战争爆发，清日围绕朝鲜和区域秩序展开海陆战。核心取证为《清德宗实录》光绪二十年相关条；军机处档、战事方略及地方奏报。编年记录适于钉住事件在朝廷文书中的时间位置，却需以地方、对方或物质材料检验其范围和后果。账本所列条件为：驻军竞争与外交谈判破裂升级为战争。其后果被限定为：战场损失、指挥责任与民众经验须多源核验。取证保留：以光绪二十年（1894年）的同期文书为定位起点；官修记录、奏报或外交文本服务特定机构立场，具体日期、规模、地方执行与对手视角须以独立材料复核。
\hypertarget{ledger:EQ-1894-YALU-NAVAL-BATTLE}{}\paragraph{光绪二十年（1894年）}黄海海战后北洋舰队受创，海军力量对比进一步恶化。核心取证为《清德宗实录》光绪二十年相关条；军机处档、战事方略及地方奏报。编年记录适于钉住事件在朝廷文书中的时间位置，却需以地方、对方或物质材料检验其范围和后果。账本所列条件为：舰艇性能、训练、战术和指挥因素共同作用。其后果被限定为：单一“装备落后”解释无法涵盖战役复杂性。取证保留：以光绪二十年（1894年）的同期文书为定位起点；官修记录、奏报或外交文本服务特定机构立场，具体日期、规模、地方执行与对手视角须以独立材料复核。
\hypertarget{ledger:EQ-1895-SHIMONOSEKI-TREATY}{}\paragraph{光绪二十一年（1895年）}《马关条约》签订，清割让台湾、澎湖并承认朝鲜独立，赔款与通商条件加重。核心取证为《清德宗实录》光绪二十一年相关条；《筹办夷务始末》相关年卷与中外照会、条约文本。这类上谕、题本或部议首先证明机构处理了何种命令、呈报与责任，而不自动证明地方收到、执行或接受。账本所列条件为：清军战败和日本军事压力迫使议和。其后果被限定为：条约的签署、割让交接和地方抵抗是不同过程。取证保留：以光绪二十一年（1895年）的同期文书为定位起点；官修记录、奏报或外交文本服务特定机构立场，具体日期、规模、地方执行与对手视角须以独立材料复核。
\hypertarget{ledger:EQ-1895-TAIWAN-REPUBLIC-RESISTANCE}{}\paragraph{光绪二十一年（1895年）}台湾割让后地方力量建立台湾民主国并抵抗日军登陆。核心取证为《清德宗实录》光绪二十一年相关条；军机处档、战事方略及地方奏报。编年记录适于钉住事件在朝廷文书中的时间位置，却需以地方、对方或物质材料检验其范围和后果。账本所列条件为：条约割让与地方社会的政治、军事反应不一致。其后果被限定为：抵抗组织与普通居民处境不能由单一民族叙事概括。取证保留：以光绪二十一年（1895年）的同期文书为定位起点；官修记录、奏报或外交文本服务特定机构立场，具体日期、规模、地方执行与对手视角须以独立材料复核。
\hypertarget{ledger:EQ-1895-TRIPLE-INTERVENTION}{}\paragraph{光绪二十一年（1895年）}三国干涉还辽改变战后列强竞争和清廷外交处境。核心取证为《清德宗实录》光绪二十一年相关条；《筹办夷务始末》相关年卷与中外照会、条约文本。这类上谕、题本或部议首先证明机构处理了何种命令、呈报与责任，而不自动证明地方收到、执行或接受。账本所列条件为：列强在东北和东亚的利益冲突加剧。其后果被限定为：干涉并非出于中立调停，需置于帝国竞争中理解。取证保留：以光绪二十一年（1895年）的同期文书为定位起点；官修记录、奏报或外交文本服务特定机构立场，具体日期、规模、地方执行与对手视角须以独立材料复核。
\hypertarget{ledger:EQ-1896-LI-LOBANOV-TREATY}{}\paragraph{光绪二十二年（1896年）}《中俄密约》及东清铁路安排加强俄国在东北的影响。核心取证为《清德宗实录》光绪二十二年相关条；《筹办夷务始末》相关年卷与中外照会、条约文本。这类上谕、题本或部议首先证明机构处理了何种命令、呈报与责任，而不自动证明地方收到、执行或接受。账本所列条件为：清廷寻求对日后的外交支撑，俄国谋求远东利益。其后果被限定为：秘密条约的文本、执行与地方影响需分层说明。取证保留：以光绪二十二年（1896年）的同期文书为定位起点；官修记录、奏报或外交文本服务特定机构立场，具体日期、规模、地方执行与对手视角须以独立材料复核。
\hypertarget{ledger:EQ-1896-MODERN-EDUCATION-EXPANSION}{}\paragraph{光绪二十二年（1896年）}甲午战后新式学堂、译书与留学讨论加速。核心取证为《清德宗实录》光绪二十二年相关条；内阁档、文集或报刊相关记录。编年记录适于钉住事件在朝廷文书中的时间位置，却需以地方、对方或物质材料检验其范围和后果。账本所列条件为：战败促使士大夫和官员反思军政、教育与技术。其后果被限定为：教育项目的地域覆盖和入学规模不宜夸大。取证保留：以光绪二十二年（1896年）的同期文书为定位起点；官修记录、奏报或外交文本服务特定机构立场，具体日期、规模、地方执行与对手视角须以独立材料复核。
\hypertarget{ledger:EQ-1897-JIAOZHOU-BAY-OCCUPATION}{}\paragraph{光绪二十三年（1897年）}德国以巨野教案为由占领胶州湾，清廷面临新的租借压力。核心取证为《清德宗实录》光绪二十三年相关条；《筹办夷务始末》相关年卷与中外照会、条约文本。这类上谕、题本或部议首先证明机构处理了何种命令、呈报与责任，而不自动证明地方收到、执行或接受。账本所列条件为：列强借教案和军舰扩张在华据点。其后果被限定为：教案、军事占领和租借谈判应分别记录。取证保留：以光绪二十三年（1897年）的同期文书为定位起点；官修记录、奏报或外交文本服务特定机构立场，具体日期、规模、地方执行与对手视角须以独立材料复核。
\hypertarget{ledger:EQ-1898-LEASED-TERRITORIES}{}\paragraph{光绪二十四年（1898年）}俄、英、法、德相继取得租借地或势力范围安排。核心取证为《清德宗实录》光绪二十四年相关条；《筹办夷务始末》相关年卷与中外照会、条约文本。这类上谕、题本或部议首先证明机构处理了何种命令、呈报与责任，而不自动证明地方收到、执行或接受。账本所列条件为：甲午后清廷军事弱势与列强竞争叠加。其后果被限定为：“势力范围”并不自动等于完整主权转移。取证保留：以光绪二十四年（1898年）的同期文书为定位起点；官修记录、奏报或外交文本服务特定机构立场，具体日期、规模、地方执行与对手视角须以独立材料复核。
\hypertarget{ledger:EQ-1898-HUNDRED-DAYS-REFORM}{}\paragraph{光绪二十四年（1898年）}光绪帝颁布变法诏令，推动官制、教育、经济与军事改革。核心取证为《清德宗实录》光绪二十四年相关条；宫中朱批奏折、内阁大库题本与《清史稿》相关志传。这类上谕、题本或部议首先证明机构处理了何种命令、呈报与责任，而不自动证明地方收到、执行或接受。账本所列条件为：甲午战败和维新思潮使改革诉求集中爆发。其后果被限定为：诏令与实际实施之间时间短、阻力大，须区别方案与结果。取证保留：以光绪二十四年（1898年）的同期文书为定位起点；官修记录、奏报或外交文本服务特定机构立场，具体日期、规模、地方执行与对手视角须以独立材料复核。
\hypertarget{ledger:EQ-1898-WUXU-COUP}{}\paragraph{光绪二十四年（1898年）}慈禧发动政变，光绪帝被幽禁，戊戌变法主要措施被撤销或改写。核心取证为《清德宗实录》光绪二十四年相关条；宫中朱批奏折、内阁大库题本与《清史稿》相关志传。这类上谕、题本或部议首先证明机构处理了何种命令、呈报与责任，而不自动证明地方收到、执行或接受。账本所列条件为：改革触及官僚利益、皇权关系与军事权力。其后果被限定为：政变叙述中的人物动机和密谋细节需谨慎处理。取证保留：以光绪二十四年（1898年）的同期文书为定位起点；官修记录、奏报或外交文本服务特定机构立场，具体日期、规模、地方执行与对手视角须以独立材料复核。
\hypertarget{ledger:EQ-1899-BOXER-UPRISING}{}\paragraph{光绪二十五年（1899年）}义和团在华北扩展，反教、反洋冲突加剧。核心取证为《清德宗实录》光绪二十五年相关条；地方志、督抚奏折及地方档案。编年记录适于钉住事件在朝廷文书中的时间位置，却需以地方、对方或物质材料检验其范围和后果。账本所列条件为：灾荒、地方武术结社、教案和列强压力相互作用。其后果被限定为：“义和团”内部成分复杂，不能以单一意识形态概括。取证保留：以光绪二十五年（1899年）的同期文书为定位起点；官修记录、奏报或外交文本服务特定机构立场，具体日期、规模、地方执行与对手视角须以独立材料复核。
\subsection{1900—1909：事件、文书与地方后果}
这一组节点按账本的日期顺序排列。相邻事件并不必然构成单一因果，却共同显示信息、资源与地方行动如何在同一时段交错。
\hypertarget{ledger:EQ-1900-BOXER-WAR}{}\paragraph{光绪二十六年（1900年）}慈禧对列强宣战，八国联军入侵并占领北京。核心取证为《清德宗实录》光绪二十六年相关条；军机处档、战事方略及地方奏报。编年记录适于钉住事件在朝廷文书中的时间位置，却需以地方、对方或物质材料检验其范围和后果。账本所列条件为：朝廷内部政策摇摆、义和团动员与外交危机相互激化。其后果被限定为：战争责任、暴力与劫掠需同时纳入中外和地方材料。取证保留：以光绪二十六年（1900年）的同期文书为定位起点；官修记录、奏报或外交文本服务特定机构立场，具体日期、规模、地方执行与对手视角须以独立材料复核。
\hypertarget{ledger:EQ-1900-SOUTHEAST-MUTUAL-PROTECTION}{}\paragraph{光绪二十六年（1900年）}东南互保使部分督抚与列强协商维持地方秩序。核心取证为《清德宗实录》光绪二十六年相关条；宫中朱批奏折、内阁大库题本与《清史稿》相关志传。这类上谕、题本或部议首先证明机构处理了何种命令、呈报与责任，而不自动证明地方收到、执行或接受。账本所列条件为：中央战争命令与地方对商贸、治安的现实考量冲突。其后果被限定为：“互保”并非整齐统一的地方独立，须逐省考察。取证保留：以光绪二十六年（1900年）的同期文书为定位起点；官修记录、奏报或外交文本服务特定机构立场，具体日期、规模、地方执行与对手视角须以独立材料复核。
\hypertarget{ledger:EQ-1901-BOXER-PROTOCOL}{}\paragraph{光绪二十七年（1901年）}《辛丑条约》签订，赔款、驻军和外交条件进一步限制清廷。核心取证为《清德宗实录》光绪二十七年相关条；《筹办夷务始末》相关年卷与中外照会、条约文本。这类上谕、题本或部议首先证明机构处理了何种命令、呈报与责任，而不自动证明地方收到、执行或接受。账本所列条件为：北京陷落和议和谈判使清廷处于弱势。其后果被限定为：赔款总额、支付机制与地方税收后果应分开分析。取证保留：以光绪二十七年（1901年）的同期文书为定位起点；官修记录、奏报或外交文本服务特定机构立场，具体日期、规模、地方执行与对手视角须以独立材料复核。
\hypertarget{ledger:EQ-1901-ZONGLI-YAMEN-REPLACED}{}\paragraph{光绪二十七年（1901年）}总理衙门改为外务部，晚清外交机构进一步制度化。核心取证为《清德宗实录》光绪二十七年相关条；宫中朱批奏折、内阁大库题本与《清史稿》相关志传。这类上谕、题本或部议首先证明机构处理了何种命令、呈报与责任，而不自动证明地方收到、执行或接受。账本所列条件为：战后改革要求重组对外事务处理。其后果被限定为：机构改名不等于外交权力和信息网络立即现代化。取证保留：以光绪二十七年（1901年）的同期文书为定位起点；官修记录、奏报或外交文本服务特定机构立场，具体日期、规模、地方执行与对手视角须以独立材料复核。
\hypertarget{ledger:EQ-1902-CIXI-RETURN-BEIJING}{}\paragraph{光绪二十八年（1902年）}慈禧与光绪帝返回北京，战后朝廷开始推出新政措施。核心取证为《清德宗实录》光绪二十八年相关条；宫中朱批奏折、内阁大库题本与《清史稿》相关志传。这类上谕、题本或部议首先证明机构处理了何种命令、呈报与责任，而不自动证明地方收到、执行或接受。账本所列条件为：辛丑议和结束流亡状态。其后果被限定为：新政既回应危机，也服务朝廷重建合法性。取证保留：以光绪二十八年（1902年）的同期文书为定位起点；官修记录、奏报或外交文本服务特定机构立场，具体日期、规模、地方执行与对手视角须以独立材料复核。
\hypertarget{ledger:EQ-1902-NEW-POLICIES-EDUCATION}{}\paragraph{光绪二十八年（1902年）}清末新政在学制、军制、商政和地方行政方面逐步启动。核心取证为《清德宗实录》光绪二十八年相关条；内阁档、文集或报刊相关记录。编年记录适于钉住事件在朝廷文书中的时间位置，却需以地方、对方或物质材料检验其范围和后果。账本所列条件为：战败与财政外交压力显示旧制难以维持。其后果被限定为：改革法令、试办与地方成效必须分别记录。取证保留：以光绪二十八年（1902年）的同期文书为定位起点；官修记录、奏报或外交文本服务特定机构立场，具体日期、规模、地方执行与对手视角须以独立材料复核。
\hypertarget{ledger:EQ-1903-BANNER-POSTS-REFORM}{}\paragraph{光绪二十九年（1903年）}清廷废除部分旗人官职垄断，调整旗籍制度与职务资格。核心取证为《清德宗实录》光绪二十九年相关条；地方志、督抚奏折及地方档案。编年记录适于钉住事件在朝廷文书中的时间位置，却需以地方、对方或物质材料检验其范围和后果。账本所列条件为：八旗财政危机与官僚改革需求上升。其后果被限定为：制度变化对不同旗籍人群的影响不均。取证保留：以光绪二十九年（1903年）的同期文书为定位起点；官修记录、奏报或外交文本服务特定机构立场，具体日期、规模、地方执行与对手视角须以独立材料复核。
\hypertarget{ledger:EQ-1903-SHANGHAI-COMMERCIAL-LAW}{}\paragraph{光绪二十九年（1903年）}商法、公司与工商管理的近代制度讨论在口岸和中央之间推进。核心取证为《清德宗实录》光绪二十九年相关条；户部奏销、盐政、河工或海关档案。这类上谕、题本或部议首先证明机构处理了何种命令、呈报与责任，而不自动证明地方收到、执行或接受。账本所列条件为：国内商业、外贸和新式企业需要新的法律语言。其后果被限定为：法典公布不等于商事审判和公司治理普遍落实。取证保留：以光绪二十九年（1903年）的同期文书为定位起点；官修记录、奏报或外交文本服务特定机构立场，具体日期、规模、地方执行与对手视角须以独立材料复核。
\hypertarget{ledger:EQ-1904-RUSSO-JAPANESE-WAR}{}\paragraph{光绪三十年（1904年）}日俄战争在中国东北爆发，清廷名义中立而东北遭受战场化。核心取证为《清德宗实录》光绪三十年相关条；军机处档、战事方略及地方奏报。编年记录适于钉住事件在朝廷文书中的时间位置，却需以地方、对方或物质材料检验其范围和后果。账本所列条件为：俄日争夺满洲铁路、港口和势力范围。其后果被限定为：中国地方社会的损失与清廷的有限行动不能被列强叙事遮蔽。取证保留：以光绪三十年（1904年）的同期文书为定位起点；官修记录、奏报或外交文本服务特定机构立场，具体日期、规模、地方执行与对手视角须以独立材料复核。
\hypertarget{ledger:EQ-1904-IMPERIAL-EXAMINATION-DEBATE}{}\paragraph{光绪三十年（1904年）}科举存废与新式教育的争论加速。核心取证为《清德宗实录》光绪三十年相关条；内阁档、文集或报刊相关记录。编年记录适于钉住事件在朝廷文书中的时间位置，却需以地方、对方或物质材料检验其范围和后果。账本所列条件为：军政危机使人才选拔制度受到质疑。其后果被限定为：讨论、诏令和实际学堂替代之间存在时间差。取证保留：以光绪三十年（1904年）的同期文书为定位起点；官修记录、奏报或外交文本服务特定机构立场，具体日期、规模、地方执行与对手视角须以独立材料复核。
\hypertarget{ledger:EQ-1905-ABOLITION-OF-CIVIL-SERVICE-EXAM}{}\paragraph{光绪三十一年（1905年）}清廷废除科举，延续千余年的取士制度终结。核心取证为《清德宗实录》光绪三十一年相关条；内阁档、文集或报刊相关记录。编年记录适于钉住事件在朝廷文书中的时间位置，却需以地方、对方或物质材料检验其范围和后果。账本所列条件为：新式教育、官僚改革与科举批评共同推动决策。其后果被限定为：废科举不等于地方教育资源立即均衡或士人身份消失。取证保留：以光绪三十一年（1905年）的同期文书为定位起点；官修记录、奏报或外交文本服务特定机构立场，具体日期、规模、地方执行与对手视角须以独立材料复核。
\hypertarget{ledger:EQ-1905-TONGMENGHUI}{}\paragraph{光绪三十一年（1905年）}中国同盟会在东京成立，革命组织和宣传网络进一步整合。核心取证为《清德宗实录》光绪三十一年相关条；地方志、督抚奏折及地方档案。编年记录适于钉住事件在朝廷文书中的时间位置，却需以地方、对方或物质材料检验其范围和后果。账本所列条件为：留日学生、海外华侨和反清团体形成跨境联系。其后果被限定为：组织自称、成员规模与地方影响需要分别核验。取证保留：以光绪三十一年（1905年）的同期文书为定位起点；官修记录、奏报或外交文本服务特定机构立场，具体日期、规模、地方执行与对手视角须以独立材料复核。
\hypertarget{ledger:EQ-1906-CONSTITUTIONAL-PROMISE}{}\paragraph{光绪三十二年（1906年）}清廷宣布预备立宪，着手考察各国制度和改革官制。核心取证为《清德宗实录》光绪三十二年相关条；宫中朱批奏折、内阁大库题本与《清史稿》相关志传。这类上谕、题本或部议首先证明机构处理了何种命令、呈报与责任，而不自动证明地方收到、执行或接受。账本所列条件为：新政需要重建统治合法性并回应立宪压力。其后果被限定为：“预备”时间表和实际权力让渡之间存在重大距离。取证保留：以光绪三十二年（1906年）的同期文书为定位起点；官修记录、奏报或外交文本服务特定机构立场，具体日期、规模、地方执行与对手视角须以独立材料复核。
\hypertarget{ledger:EQ-1906-NEW-ARMY-REFORM}{}\paragraph{光绪三十二年（1906年）}新军编练与军事学堂扩展，军队成为新式政治动员的重要空间。核心取证为《清德宗实录》光绪三十二年相关条；军机处档、战事方略及地方奏报。编年记录适于钉住事件在朝廷文书中的时间位置，却需以地方、对方或物质材料检验其范围和后果。账本所列条件为：甲午、庚子和日俄战争显示旧军制弱点。其后果被限定为：编制、装备和忠诚网络在各省差异很大。取证保留：以光绪三十二年（1906年）的同期文书为定位起点；官修记录、奏报或外交文本服务特定机构立场，具体日期、规模、地方执行与对手视角须以独立材料复核。
\hypertarget{ledger:EQ-1907-MANCHURIAN-PROVINCES}{}\paragraph{光绪三十三年（1907年）}清廷将东三省调整为行省体制，加强近代行政与边防经营。核心取证为《清德宗实录》光绪三十三年相关条；军机处档、理藩院档或相关方略。编年记录适于钉住事件在朝廷文书中的时间位置，却需以地方、对方或物质材料检验其范围和后果。账本所列条件为：日俄战争后东北主权和铁路问题尖锐化。其后果被限定为：省制设立与铁路沿线列强控制并存。取证保留：以光绪三十三年（1907年）的同期文书为定位起点；官修记录、奏报或外交文本服务特定机构立场，具体日期、规模、地方执行与对手视角须以独立材料复核。
\hypertarget{ledger:EQ-1907-QING-LEGAL-REFORM}{}\paragraph{光绪三十三年（1907年）}清廷推进修律、审判和警政改革，尝试改造传统法律制度。核心取证为《清德宗实录》光绪三十三年相关条；宫中朱批奏折、内阁大库题本与《清史稿》相关志传。这类上谕、题本或部议首先证明机构处理了何种命令、呈报与责任，而不自动证明地方收到、执行或接受。账本所列条件为：领事裁判权和国内治理压力推动法制改革。其后果被限定为：新法文本、试行和基层司法实践不能混同。取证保留：以光绪三十三年（1907年）的同期文书为定位起点；官修记录、奏报或外交文本服务特定机构立场，具体日期、规模、地方执行与对手视角须以独立材料复核。
\hypertarget{ledger:EQ-1908-GUANGXU-CIXI-DEATHS}{}\paragraph{光绪三十四年（1908年）}光绪帝与慈禧太后相继去世，溥仪继位为宣统帝。核心取证为《清德宗实录》光绪三十四年相关条；宫中朱批奏折、内阁大库题本与《清史稿》相关志传。这类上谕、题本或部议首先证明机构处理了何种命令、呈报与责任，而不自动证明地方收到、执行或接受。账本所列条件为：晚清权力核心突然更替。其后果被限定为：摄政王载沣与幼帝体制面临改革和危机双重压力。取证保留：以光绪三十四年（1908年）的同期文书为定位起点；官修记录、奏报或外交文本服务特定机构立场，具体日期、规模、地方执行与对手视角须以独立材料复核。
\hypertarget{ledger:EQ-1908-CONSTITUTIONAL-OUTLINE}{}\paragraph{光绪三十四年（1908年）}清廷公布《钦定宪法大纲》等文件，明确预备立宪框架。核心取证为《清德宗实录》光绪三十四年相关条；宫中朱批奏折、内阁大库题本与《清史稿》相关志传。这类上谕、题本或部议首先证明机构处理了何种命令、呈报与责任，而不自动证明地方收到、执行或接受。账本所列条件为：立宪压力与中央保留权力的意图并存。其后果被限定为：文本承诺与咨议局、资政院实际权限存在落差。取证保留：以光绪三十四年（1908年）的同期文书为定位起点；官修记录、奏报或外交文本服务特定机构立场，具体日期、规模、地方执行与对手视角须以独立材料复核。
